% !PRINTED 286
\weberpage{286}
Unter den Theilern dieser Gruppe sind hervorzuheben die
Octa\"edergruppe
\[
\chi^\lambda \omega^u \Theta^v
\]
vom Index 7, sodann die Elemente
\[
\tau^\varrho \chi^\lambda,
\]
die nach den Relationen (12) und (13) eine Gruppe 21\textsuperscript{sten} Grades,
also einen Theiler der Gesammtgruppe vom Index 8 bilden;
ferner die Gruppe 8\textsuperscript{ten} Grades $\omega^u \Theta^v$. Die gesammte Gruppe
lässt sich auch in anderer Reihenfolge so darstellen:
\[
\chi^\lambda \omega^u \Theta^v \tau^\varrho, \quad \omega^u \Theta^v \chi^\lambda \tau^\varrho, \quad \tau^\varrho \omega^u \Theta^v \chi^\lambda.
\]

\bigskip
\begin{center}
\rule{4cm}{0.4pt}
\end{center}
